%=============================================================================
% Engineering Design Document:
% 4-Node ψ-Linked Quantum Sensor Array Demonstrator
%
% Patent #6 — Density Field Dynamics
% Author: Gary Thomas Alcock
%=============================================================================

\documentclass[12pt,a4paper]{article}

\usepackage{amsmath,amssymb}
\usepackage{physics}
\usepackage{graphicx}
\usepackage{geometry}
\usepackage{hyperref}
\usepackage{booktabs}
\usepackage{float}
\usepackage{enumitem}
\usepackage{tcolorbox}
\usepackage{longtable}
\usepackage{multirow}
\usepackage{array}

\geometry{margin=2.5cm}

\title{\textbf{Engineering Design Document:\\
4-Node $\psi$-Linked Quantum Sensor Array Demonstrator}\\[6pt]
\large Patent \#6 Supporting Engineering Specification\\
\large Density Field Dynamics Research Programme}

\author{Gary Thomas Alcock\\
\texttt{gary@densityfielddynamics.com}}

\date{March 2026 \quad Rev.\ 1.0}

\begin{document}

\maketitle

\begin{abstract}
This document provides the complete engineering design for a 4-node
$\psi$-linked quantum sensor array demonstrator. The system uses
$^{87}$Rb atom interferometers as the primary $\psi$-field sensors,
arranged in a tetrahedral configuration with baselines from 10 m to
10 km, to verify the DFD prediction that distributed sensors sharing
a common gravitational $\psi$ envelope exhibit phase correlations
beyond classical environmental coupling. The design is specified to
sufficient detail for fabrication and deployment, including sensor
node architecture, environmental monitoring subsystems, data
acquisition pipeline, systematic error budget, calibration procedures,
and a phased deployment plan.
\end{abstract}

\tableofcontents
\newpage

%=============================================================================
\section{System Requirements}\label{sec:requirements}
%=============================================================================

\subsection{Primary Requirements}

\begin{enumerate}[label=\textbf{PR-\arabic*}]
\item Measure inter-node phase correlations with precision
$\delta C_{ij} < 10^{-7}$ after $10^6$ measurement cycles.
\item Reject all known classical correlation channels (seismic,
electromagnetic, thermal, tidal) to the level of $10^{-8}$ in the
correlation coefficient.
\item Operate with baselines $d = 10$ m to 10 km without physical
coherent links between nodes.
\item Achieve single-shot phase sensitivity $< 10^{-2}$ rad per node.
\item Provide full 3D measurement capability through non-planar
(tetrahedral) node placement.
\end{enumerate}

\subsection{Derived Requirements}

\begin{enumerate}[label=\textbf{DR-\arabic*}]
\item Each node shall contain an independent $^{87}$Rb cold-atom
interferometer with interrogation time $T = 1$--$10$ s.
\item Inter-node timing synchronization $< 10$ ns via GPS common-view.
\item Environmental monitoring at each node with $\geq 6$ independent
channels.
\item Data recording at each node with absolute timestamp, stored
locally and transmitted via standard network.
\item The system shall be operable in three deployment phases
(compact, extended, continental) without node redesign.
\end{enumerate}

\subsection{Performance Targets}

\begin{table}[H]
\centering
\begin{tabular}{lcc}
\toprule
\textbf{Metric} & \textbf{Phase 1} & \textbf{Phase 2/3} \\
\midrule
Baseline & 10--100 m & 1--10$^4$ km \\
Phase sensitivity/shot & $10^{-2}$ rad & $10^{-3}$ rad \\
Correlation precision ($10^6$ shots) & $10^{-5}$ & $10^{-7}$ \\
Environmental rejection & $10^{-6}$ & $10^{-8}$ \\
Measurement cadence & 0.1 Hz & 0.1 Hz \\
Continuous operation & $>$ 30 days & $>$ 90 days \\
\bottomrule
\end{tabular}
\caption{Performance targets by deployment phase.}
\label{tab:performance}
\end{table}


%=============================================================================
\section{System Architecture}\label{sec:architecture}
%=============================================================================

\subsection{Top-Level Block Diagram}

The system consists of four identical, autonomous sensor nodes plus a
central data processing facility:

\begin{tcolorbox}[colback=gray!5, colframe=gray!60, title=System Architecture]
\begin{verbatim}
+------------------+     +------------------+
|   NODE A         |     |   NODE B         |
|  Atom IF + Env   |     |  Atom IF + Env   |
|  Local Clock     |     |  Local Clock     |
|  Local DAQ       |     |  Local DAQ       |
+--------+---------+     +--------+---------+
         |                         |
         |    Standard Network     |
         +------------+------------+
                      |
              +-------+--------+
              | CENTRAL        |
              | PROCESSING     |
              | Correlation    |
              | Analysis       |
              +-------+--------+
                      |
         +------------+------------+
         |                         |
+--------+---------+     +--------+---------+
|   NODE C         |     |   NODE D         |
|  Atom IF + Env   |     |  Atom IF + Env   |
|  Local Clock     |     |  Local Clock     |
|  Local DAQ       |     |  Local DAQ       |
+------------------+     +------------------+
\end{verbatim}
\end{tcolorbox}

\textbf{Key design principle}: Each node is \emph{fully autonomous}.
There is no real-time coherent link between nodes. The $\psi$ field
itself provides the correlation; the engineering challenge is to
measure it without introducing spurious correlations through the
hardware.

\subsection{Node Functional Decomposition}

Each node contains six subsystems:

\begin{enumerate}
\item \textbf{Atom Interferometer Core} (AIC): Cold-atom preparation,
Raman interrogation, state detection
\item \textbf{Laser Subsystem} (LAS): Cooling, repumping, Raman, and
reference lasers
\item \textbf{Environmental Monitoring Suite} (EMS): Seismic, magnetic,
thermal, barometric, tilt sensors
\item \textbf{Timing and Synchronization} (TAS): Atomic clock, GPS
receiver, timestamp generator
\item \textbf{Data Acquisition and Control} (DAC): Digitization,
local storage, sequence control
\item \textbf{Infrastructure} (INF): Vacuum, magnetic shielding, thermal
enclosure, power, network
\end{enumerate}


%=============================================================================
\section{Atom Interferometer Core}\label{sec:aic}
%=============================================================================

\subsection{Atom Source}

\begin{table}[H]
\centering
\begin{tabular}{ll}
\toprule
\textbf{Parameter} & \textbf{Specification} \\
\midrule
Species & $^{87}$Rb \\
Cooling transition & $5S_{1/2}(F=2) \to 5P_{3/2}(F'=3)$, 780.241 nm \\
MOT loading rate & $\geq 10^9$ atoms/s \\
MOT atom number & $\geq 10^9$ atoms \\
Optical molasses temperature & $\leq 3\,\mu$K \\
Evaporative cooling (optional) & $\leq 50$ nK (BEC regime) \\
Launch velocity (fountain) & 0--5 m/s (vertical) \\
Atom number at interrogation & $\geq 10^6$ atoms \\
State preparation & $|F=1, m_F=0\rangle$ via microwave selection \\
\bottomrule
\end{tabular}
\caption{Atom source specifications.}
\end{table}

\subsection{Interferometer Geometry}

\textbf{Configuration}: Vertical Mach-Zehnder with three-pulse sequence
$\pi/2$--$\pi$--$\pi/2$ at times $t = 0, T, 2T$.

\textbf{Raman transitions}:
\begin{itemize}
\item Two-photon Raman via $5P_{3/2}$
\item Counter-propagating beams: $k_{\text{eff}} = 2k = 1.61 \times 10^7$ m$^{-1}$
\item Recoil velocity: $v_r = \hbar k_{\text{eff}}/m = 1.18 \times 10^{-2}$ m/s
\item Rabi frequency: $\Omega_R \geq 2\pi \times 25$ kHz ($\pi$-pulse
duration $\leq 20\,\mu$s)
\end{itemize}

\textbf{Interrogation time}: Adjustable $T = 0.1$--$10$ s.

\textbf{Phase shift (gravitational)}:
\begin{equation}
\phi_{\text{grav}} = k_{\text{eff}}\,g\,T^2.
\end{equation}
For $T = 1$ s: $\phi_{\text{grav}} = 1.58 \times 10^8$ rad.

\textbf{DFD-specific phase}:
\begin{equation}
\phi_{\text{DFD}} = \frac{\hbar k_{\text{eff}}^2}{m}\frac{g}{c^2}T^3.
\end{equation}
For $T = 1$ s: $\phi_{\text{DFD}} \approx 2 \times 10^{-11}$ rad.

\subsection{Detection System}

\begin{itemize}
\item \textbf{Method}: Normalized fluorescence detection
\item \textbf{Probe beam}: Resonant with $F=2 \to F'=3$ cycling
transition, 3 ms pulse
\item \textbf{Detection}: CCD or CMOS camera imaging the atom cloud;
separate images for $|F=1\rangle$ and $|F=2\rangle$ populations
\item \textbf{Population ratio}:
$P = N_2/(N_1 + N_2) = \frac{1}{2}(1 + \cos\phi)$
\item \textbf{Detection noise}: Quantum projection noise limited:
$\delta\phi = 1/\sqrt{N_{\text{atom}}} \approx 10^{-3}$ rad for
$10^6$ detected atoms
\end{itemize}

\subsection{Vacuum System}

\begin{table}[H]
\centering
\begin{tabular}{ll}
\toprule
\textbf{Parameter} & \textbf{Specification} \\
\midrule
Chamber material & Titanium (low magnetic permeability) \\
Base pressure & $\leq 5 \times 10^{-11}$ mbar \\
Pumping & 40 L/s ion pump + Ti sublimation pump \\
Rb source & Alkali metal dispenser (AMD) \\
Viewport material & UV-grade fused silica, AR coated at 780 nm \\
Bakeout temperature & 200$^\circ$C for 48 h \\
Interferometry volume & $\geq 30$ cm vertical $\times$ 5 cm diameter \\
Total chamber length & 1.5 m (vertical, for $T = 1$ s fountain) \\
\bottomrule
\end{tabular}
\caption{Vacuum system specifications.}
\end{table}

For extended interrogation times ($T > 1$ s), the chamber extends
to accommodate the atomic fountain trajectory:
\begin{equation}
h_{\max} = \frac{1}{2}g T^2 + v_0 T,
\end{equation}
where $v_0$ is the launch velocity. For $T = 3$ s with zero apex
velocity: $h_{\max} = \frac{1}{2}(9.8)(3)^2 = 44$ m. This requires
a tall vacuum tower, achieved in Phase 2/3 using existing infrastructure
(e.g., elevator shafts, mine shafts, or purpose-built towers similar
to those at Stanford or Wuhan).


%=============================================================================
\section{Laser Subsystem}\label{sec:laser}
%=============================================================================

\subsection{Laser Architecture}

\begin{table}[H]
\centering
\begin{tabular}{llll}
\toprule
\textbf{Function} & \textbf{Wavelength} & \textbf{Power} & \textbf{Linewidth} \\
\midrule
Cooling/detection & 780.241 nm & 200 mW & $< 1$ MHz \\
Repumper & 780.232 nm & 50 mW & $< 1$ MHz \\
Raman 1 & 780.241 nm & 100 mW & See below \\
Raman 2 & 780.241 nm + 6.835 GHz & 100 mW & See below \\
Dipole trap (opt.) & 1064 nm & 10 W & N/A \\
\bottomrule
\end{tabular}
\caption{Laser specifications by function.}
\end{table}

\subsection{Raman Laser Phase Noise}

The Raman laser pair must have low relative phase noise, as this
directly maps onto interferometer phase noise. Requirements:

\begin{itemize}
\item \textbf{Phase noise spectral density}:
$S_\phi(f) < 10^{-4}$ rad$^2$/Hz at $f = 1$ Hz
\item \textbf{Implementation}: Phase-locked loop (PLL) locking two
ECDLs to a common ultra-stable cavity, with the 6.835 GHz offset
derived from a low-noise microwave synthesizer referenced to the local
atomic clock
\item \textbf{Fiber noise cancellation}: Active fiber noise cancellation
on all paths from laser table to vacuum chamber
\end{itemize}

\subsection{Reference Cavity}

Each node includes an ultra-stable optical cavity for the local
laser frequency reference:
\begin{itemize}
\item \textbf{Spacer}: ULE glass, 10 cm length
\item \textbf{Finesse}: $\geq 100{,}000$
\item \textbf{Thermal noise floor}: $\sim 10^{-16}$ at 1 s
\item \textbf{Temperature control}: $\pm 1$ mK at zero-CTE temperature
\item \textbf{Vibration isolation}: Notch-and-spring mount on active
platform
\end{itemize}

\textbf{Critical note}: The reference cavity provides a local frequency
standard for each node. It does \emph{not} need to be coherent with
cavities at other nodes. The $\psi$-correlation analysis uses only the
\emph{phase of the atom interferometer output}, not the laser frequency.


%=============================================================================
\section{Environmental Monitoring Suite}\label{sec:ems}
%=============================================================================

\subsection{Sensor Complement}

Each node deploys the following environmental sensors:

\begin{table}[H]
\centering
\small
\begin{tabular}{llll}
\toprule
\textbf{Channel} & \textbf{Sensor type} & \textbf{Sensitivity} & \textbf{Bandwidth} \\
\midrule
Seismic (3-axis) & Broadband seismometer & $10^{-10}$ m/s$^2$/$\sqrt{\text{Hz}}$ & 0.01--100 Hz \\
Tilt (2-axis) & Precision tiltmeter & 1 nrad/$\sqrt{\text{Hz}}$ & DC--1 Hz \\
Magnetic (3-axis) & Fluxgate magnetometer & 10 pT/$\sqrt{\text{Hz}}$ & DC--100 Hz \\
Temperature (4 pts) & Pt100 RTD & 0.1 mK & DC--0.1 Hz \\
Barometric & Microbarograph & 0.01 Pa/$\sqrt{\text{Hz}}$ & 0.001--10 Hz \\
Humidity & Capacitive sensor & 0.5\% RH & DC--0.01 Hz \\
Gravity gradient & Tidal gravimeter & $10^{-11}$ m/s$^2$/$\sqrt{\text{Hz}}$ & DC--0.01 Hz \\
EM field (RF) & Spectrum analyzer & $-120$ dBm & 1 MHz--6 GHz \\
\bottomrule
\end{tabular}
\caption{Environmental monitoring sensor complement per node.}
\label{tab:env-sensors}
\end{table}

\subsection{Environmental Rejection Strategy}

The environmental data serves two functions:
\begin{enumerate}
\item \textbf{Real-time veto}: Shots taken during large transients
(earthquakes, nearby vehicle traffic, magnetic storms) are flagged
and optionally excluded.
\item \textbf{Post-processing regression}: Environmental channels are
regressed against the interferometer phase to identify and remove
classical correlations:
\begin{equation}
\phi_i^{\text{clean}} = \phi_i^{\text{raw}} - \sum_k \beta_{ik}\,e_k(t_i),
\end{equation}
where $e_k$ is the $k$-th environmental channel and $\beta_{ik}$ are
regression coefficients determined from calibration data.
\end{enumerate}

\subsection{Witness Sensors for Spurious Correlations}

To ensure that any detected inter-node correlation is not due to
common environmental driving, the following witness channels are
specifically monitored at all nodes simultaneously:

\begin{itemize}
\item \textbf{Schumann resonances} (7.83, 14.3, 20.8, 27.3 Hz):
Earth-ionosphere cavity modes that are globally correlated; measured
via magnetic field sensor
\item \textbf{Tidal gravity} (12.42-hour and 24-hour periods): Measured
by tidal gravimeter and compared with theoretical tidal model
\item \textbf{Atmospheric pressure loading}: Measured by microbarograph;
correlated across nodes at synoptic scales
\item \textbf{Seismic surface waves}: Measured by seismometers; known
propagation velocity allows time-delay prediction between nodes
\end{itemize}


%=============================================================================
\section{Timing and Synchronization}\label{sec:timing}
%=============================================================================

\subsection{Clock Architecture}

Each node contains:
\begin{itemize}
\item \textbf{Local atomic clock}: Rubidium frequency standard
(e.g., Microsemi SA.45s CSAC or equivalent)
\begin{itemize}
\item Short-term stability: $\sigma_y(1\,\text{s}) < 3 \times 10^{-10}$
\item Medium-term: $\sigma_y(100\,\text{s}) < 3 \times 10^{-12}$
\end{itemize}
\item \textbf{GPS receiver}: Dual-frequency (L1/L2), timing mode
\begin{itemize}
\item Absolute timing: $< 20$ ns (1$\sigma$) to UTC
\item Common-view differential: $< 5$ ns between nodes
\end{itemize}
\item \textbf{Timestamp counter}: FPGA-based, 1 ns resolution, triggered
by each interferometer measurement cycle
\end{itemize}

\subsection{Synchronization Protocol}

\begin{enumerate}
\item All nodes reference GPS 1-PPS (pulse per second) for coarse
synchronization ($< 100$ ns).
\item Common-view GPS processing provides differential timing to $< 5$ ns.
\item Each interferometer shot is timestamped with the local FPGA counter
referenced to the GPS-disciplined Rb clock.
\item Post-processing applies common-view corrections to align all
timestamps to a common timescale.
\end{enumerate}

\textbf{Synchronization requirement for $\psi$-correlation}: The atom
interferometer interrogation time $T = 1$--$10$ s means that the
synchronization error must satisfy $\delta t \ll T$. Even at $\delta t
= 100$ ns, the fractional timing error is $10^{-7}$ relative to $T = 1$ s,
which is negligible for the correlation measurement.


%=============================================================================
\section{Data Acquisition and Processing}\label{sec:daq}
%=============================================================================

\subsection{Local Data Acquisition}

Each node records per shot:
\begin{itemize}
\item Interferometer phase: $\phi_i$ (derived from population ratio)
\item Atom number: $N_1, N_2$ (for normalization and noise estimation)
\item Timestamp: $t_i$ (FPGA counter, GPS-referenced)
\item Environmental data: all channels at matching cadence
\item Diagnostics: laser power, MOT fluorescence, vacuum pressure,
shield temperature
\end{itemize}

\textbf{Data rate per node}: $\sim 1$ kB/shot $\times$ 0.1 Hz $= 100$ B/s
continuous $\approx 8.6$ MB/day. Negligible bandwidth for any network link.

\subsection{Data Transfer}

\begin{itemize}
\item \textbf{Phase 1} (compact): Local Ethernet, real-time transfer
\item \textbf{Phase 2} (extended): VPN over university/commercial internet
\item \textbf{Phase 3} (continental): Same as Phase 2; latency of hours
is acceptable since no real-time coherent processing is needed
\end{itemize}

\subsection{Central Processing Pipeline}

The central processing facility performs the following analysis chain:

\begin{enumerate}
\item \textbf{Data ingestion}: Collect shot records from all four nodes,
align timestamps.

\item \textbf{Quality cuts}: Remove shots where:
\begin{itemize}
\item Atom number $< N_{\min} = 5 \times 10^5$
\item Environmental transient exceeds threshold
\item Laser diagnostics indicate unlock
\end{itemize}

\item \textbf{Environmental regression}: For each node, compute cleaned
phase:
\begin{equation}
\phi_i^{(\text{clean})} = \phi_i - \sum_{k=1}^{N_{\text{env}}} \hat{\beta}_{ik}\,e_{ik}
\end{equation}

\item \textbf{Correlation computation}: For each pair $(i,j)$:
\begin{equation}
\hat{C}_{ij} = \frac{1}{M}\sum_{m=1}^{M}\cos\!\left(\phi_i^{(\text{clean})}(m) - \phi_j^{(\text{clean})}(m)\right)
\end{equation}

\item \textbf{Null tests}:
\begin{itemize}
\item Time-reversed correlation (should give same result for $\psi$, not for causal environmental coupling)
\item Frequency-domain correlation vs.\ known environmental spectra
\item Jackknife/bootstrap error estimation
\end{itemize}

\item \textbf{Signal extraction}: Excess correlation
$\Delta C_{ij} = \hat{C}_{ij} - C_{ij}^{(\text{null})}$ tested against
the DFD prediction:
\begin{equation}
C_{ij}^{(\psi)} = \exp\!\left[-\sigma_\psi^2\!\left(1 - \frac{\Xi(d_{ij})}{\sigma_\psi^2}\right)\right]
\end{equation}
\end{enumerate}

\subsection{Statistical Framework}

\textbf{Detection threshold}: A $5\sigma$ excess correlation above the
null model, where $\sigma$ is estimated from bootstrap resampling of the
data with randomized time offsets between nodes (destroying any true
temporal correlation while preserving single-node statistics).

\textbf{Required measurement cycles}: For a target correlation precision
$\delta C \sim 10^{-7}$:
\begin{equation}
M \geq \frac{1}{\delta C^2} = 10^{14}\,\text{shots},
\end{equation}
at 0.1 Hz cadence, this requires $10^{15}$ s $\approx 3 \times 10^7$ years
if limited to single-shot precision. However, with single-shot noise
$\sigma_\phi \sim 10^{-3}$ rad and the correlation signal being a function
of the \emph{common-mode} $\psi$ variation (not the per-shot noise):
\begin{equation}
\delta\hat{C}_{ij} = \frac{\sigma_\phi^2}{\sqrt{M}} \approx \frac{10^{-6}}{\sqrt{M}}.
\end{equation}
For $\delta\hat{C} = 10^{-7}$: $M = 10^{2}$ shots $\approx 1000$ s.

The bottleneck is not statistical precision but \emph{systematic rejection}:
demonstrating that the measured correlation is not environmental requires
the extensive witness sensor and null test program described above.


%=============================================================================
\section{Magnetic Shielding}\label{sec:shielding}
%=============================================================================

\subsection{Requirements}

Magnetic fields couple to atom interferometers through the quadratic
Zeeman effect. For the $m_F = 0 \to m_F = 0$ clock transition:
\begin{equation}
\delta\phi_B = 2\pi \times 575.15\,\text{Hz/G}^2 \times B^2 \times T.
\end{equation}
For $\delta\phi_B < 10^{-4}$ rad at $T = 1$ s:
$B < \sqrt{10^{-4}/(2\pi \times 575.15)} \approx 5\,\mu$G $= 0.5$ nT.

\subsection{Shield Design}

Four-layer mu-metal shield surrounding the vacuum chamber:
\begin{itemize}
\item Layer 1 (outermost): 1.5 mm Mu-metal, shielding factor $\sim 100$
\item Layer 2: 1.0 mm Mu-metal, shielding factor $\sim 80$
\item Layer 3: 1.0 mm Mu-metal, shielding factor $\sim 80$
\item Layer 4 (innermost): 0.5 mm Mu-metal, shielding factor $\sim 50$
\item \textbf{Total shielding factor}: $\geq 10^7$ (DC)
\item \textbf{Residual field}: $< 50$ pT (Earth's field $\sim 50\,\mu$T
$/10^7$)
\end{itemize}

Internal bias coils provide a uniform 100 mG ($10\,\mu$T) quantization
field along the vertical axis.


%=============================================================================
\section{Vibration Isolation}\label{sec:vibration}
%=============================================================================

\subsection{Requirements}

The atom interferometer phase is sensitive to vibrations of the Raman
retroreflection mirror:
\begin{equation}
\delta\phi_{\text{vib}} = k_{\text{eff}}\,\delta x_{\text{mirror}}.
\end{equation}
For $\delta\phi_{\text{vib}} < 10^{-2}$ rad:
$\delta x < 6 \times 10^{-10}$ m $= 0.6$ nm.

\subsection{Isolation Strategy}

\begin{enumerate}
\item \textbf{Passive platform}: Minus-K or similar negative-stiffness
isolator
\begin{itemize}
\item Vertical resonance: 0.5 Hz
\item Isolation above 1 Hz: $> 40$ dB/decade
\end{itemize}

\item \textbf{Active feedback}: Seismometer-driven piezo actuator on
retroreflection mirror mount
\begin{itemize}
\item Bandwidth: 0.01--50 Hz
\item Residual acceleration: $< 10^{-9}$ m/s$^2$/$\sqrt{\text{Hz}}$
at 1 Hz
\end{itemize}

\item \textbf{Post-correction}: Correlate seismometer signal with
interferometer phase; subtract in post-processing
\begin{itemize}
\item Additional factor $\sim 100$ rejection at frequencies where
seismometer SNR is high
\end{itemize}
\end{enumerate}

\subsection{Differential Vibration Rejection}

For the $\psi$-correlation measurement, the critical quantity is the
\emph{differential} vibration between nodes. At separations $> 100$ m,
ground vibrations are largely uncorrelated above $\sim 0.1$ Hz (except
for coherent seismic waves). The environmental regression removes
correlated seismic signals; the uncorrelated component averages down
as $1/\sqrt{M}$.


%=============================================================================
\section{Thermal Control}\label{sec:thermal}
%=============================================================================

\subsection{Requirements}

Temperature fluctuations affect the interferometer through:
\begin{itemize}
\item Blackbody radiation shift of the clock transition:
$\partial\nu_{\text{BBR}}/\partial T \approx 0.02$ Hz/K at 300 K
\item Vacuum chamber outgassing rate (affects collisional phase shift)
\item Reference cavity frequency drift
\end{itemize}

\subsection{Thermal Enclosure}

\begin{itemize}
\item \textbf{Outer insulation}: 10 cm expanded polystyrene, covering
the full apparatus
\item \textbf{Active temperature control}: PID-controlled heater tape
on the vacuum chamber, stability $\pm 1$ mK over 24 hours
\item \textbf{Reference cavity}: Independent thermal enclosure with
$\pm 0.1$ mK stability near the ULE zero-CTE temperature ($\sim 30$--$35^\circ$C)
\item \textbf{Air conditioning}: Laboratory HVAC maintaining $\pm 0.5$ K
ambient temperature
\end{itemize}


%=============================================================================
\section{Node Physical Layout}\label{sec:physical}
%=============================================================================

\subsection{Compact Node (Phase 1)}

Dimensions: 2.0 m (H) $\times$ 1.0 m (W) $\times$ 0.8 m (D)
\begin{itemize}
\item Lower section (0--0.5 m): Electronics, DAQ, power supplies
\item Middle section (0.5--1.5 m): Vacuum chamber (vertical), magnetic
shield, vibration isolation platform
\item Upper section (1.5--2.0 m): Optical breadboard with laser system
\item Side panel: Environmental sensor cluster
\end{itemize}

Total mass per node: $\sim 500$ kg (including shielding).

\subsection{Extended Node (Phase 2/3)}

For $T > 1$ s, the vacuum tower extends above the compact node:
\begin{itemize}
\item Tower height: 5 m ($T = 1$ s) to 50 m ($T = 3$ s)
\item Tower diameter: 30 cm (vacuum tube) within 1 m (thermal/magnetic
shield)
\item Foundation: Vibration-isolated pier, separate from building structure
\end{itemize}


%=============================================================================
\section{Calibration Procedures}\label{sec:calibration}
%=============================================================================

\subsection{Phase Calibration}

\begin{enumerate}
\item \textbf{Gravity measurement}: Each node independently measures
$g$ from the standard $k_{\text{eff}}gT^2$ phase:
\begin{equation}
g_i = \frac{\phi_i}{k_{\text{eff}}T^2}
\end{equation}
This calibrates the absolute phase-to-acceleration conversion.

\item \textbf{Raman laser calibration}: Modulate the Raman phase at
known amplitude and frequency; verify the phase response matches.

\item \textbf{Tidal calibration}: Compare the measured tidal phase
modulation at each node with the theoretical solid-Earth tide:
\begin{equation}
\delta g_{\text{tide}}(t) = \sum_{n} A_n \cos(\omega_n t + \phi_n),
\end{equation}
where the principal constituents ($M_2, S_2, K_1, O_1$, etc.) have
known amplitudes and phases. This provides an absolute calibration of
the $\psi$-dependent phase at the $\sim 10^{-7}$ level.
\end{enumerate}

\subsection{Correlation Calibration}

\begin{enumerate}
\item \textbf{Self-correlation}: Each node's autocorrelation $C_{ii}$
should be $\leq 1$, decreasing with the Allan variance of the
measurement noise.

\item \textbf{Known-signal injection}: Introduce a known gravitational
signal (e.g., a moving test mass near one node) and verify:
\begin{itemize}
\item The signal appears in the near-node data
\item It does \emph{not} appear in the far-node data (classical test)
\item It does appear in the $\psi$-correlation at the predicted level
(DFD test)
\end{itemize}

\item \textbf{Hardware exchange}: Swap physical sensor nodes between
positions to verify that any detected correlation follows the
\emph{position}, not the \emph{hardware}.
\end{enumerate}


%=============================================================================
\section{Deployment Plan}\label{sec:deployment}
%=============================================================================

\subsection{Phase 1: Compact Demonstrator (Year 1--2)}

\begin{table}[H]
\centering
\begin{tabular}{lp{10cm}}
\toprule
\textbf{Month} & \textbf{Activity} \\
\midrule
1--6 & Design and procure four identical atom interferometer nodes.
Assemble laser systems, vacuum chambers, magnetic shields. \\
7--9 & Commission each node individually: achieve cold-atom loading,
Raman interferometry, gravity measurement to $10^{-8}$ m/s$^2$. \\
10--12 & Deploy four nodes in tetrahedral configuration within
laboratory campus (baselines 10--100 m). Commission environmental
monitoring. \\
13--18 & Collect correlation data ($\geq 10^6$ shots). Perform null
tests and environmental regression. \\
19--24 & Analysis, systematic studies, hardware exchange tests.
Publish Phase 1 results. \\
\bottomrule
\end{tabular}
\end{table}

\subsection{Phase 2: Metropolitan Scale (Year 2--4)}

\begin{table}[H]
\centering
\begin{tabular}{lp{10cm}}
\toprule
\textbf{Month} & \textbf{Activity} \\
\midrule
25--30 & Relocate nodes to university/laboratory sites across a
metropolitan area (baselines 1--10 km). Upgrade vacuum towers for
$T \leq 3$ s. \\
31--36 & Commission at new sites. Establish GPS common-view timing.
Verify independent operation. \\
37--48 & Extended data collection ($\geq 10^7$ shots over 1 year).
Sensitivity to gravitational wave signals in the 0.01--1 Hz band.
$\psi$-field tomography of local geology. \\
\bottomrule
\end{tabular}
\end{table}

\subsection{Phase 3: Continental Scale (Year 4--7)}

\begin{table}[H]
\centering
\begin{tabular}{lp{10cm}}
\toprule
\textbf{Month} & \textbf{Activity} \\
\midrule
49--60 & Partner with existing atom interferometry laboratories
(e.g., MAGIS, MIGA, ZAIGA sites). Install $\psi$-correlation
DAQ at partner nodes. \\
61--72 & Continental-baseline correlation analysis. Gravitational
wave searches. Earth interior tomography. \\
73--84 & Full 4-node continental network science operations.
Publish definitive $\psi$-correlation results. \\
\bottomrule
\end{tabular}
\end{table}


%=============================================================================
\section{Cost Estimate}\label{sec:cost}
%=============================================================================

\begin{table}[H]
\centering
\begin{tabular}{lrr}
\toprule
\textbf{Item} & \textbf{Per node} & \textbf{4 nodes} \\
\midrule
Vacuum system (chamber, pumps, gauges) & \$120k & \$480k \\
Laser system (cooling, Raman, reference) & \$200k & \$800k \\
Optics (AOM, EOM, fibers, mirrors) & \$80k & \$320k \\
Magnetic shielding (4-layer) & \$40k & \$160k \\
Vibration isolation platform & \$30k & \$120k \\
Environmental sensors suite & \$50k & \$200k \\
DAQ electronics and FPGA & \$30k & \$120k \\
GPS timing receiver & \$15k & \$60k \\
Thermal enclosure & \$10k & \$40k \\
\midrule
\textbf{Node subtotal} & \textbf{\$575k} & \textbf{\$2.30M} \\
\midrule
Central processing facility & --- & \$100k \\
Installation and commissioning & --- & \$200k \\
Contingency (20\%) & --- & \$520k \\
\midrule
\textbf{Phase 1 total} & --- & \textbf{\$3.12M} \\
\bottomrule
\end{tabular}
\caption{Phase 1 cost estimate (compact demonstrator).}
\label{tab:cost}
\end{table}

Phase 2 additional costs: $\sim\$1$M (tower construction, relocation,
extended operation).

Phase 3 additional costs: $\sim\$0.5$M (partner integration, travel,
data infrastructure).


%=============================================================================
\section{Risk Assessment}\label{sec:risk}
%=============================================================================

\begin{table}[H]
\centering
\small
\begin{tabular}{p{4cm}ccp{5cm}}
\toprule
\textbf{Risk} & \textbf{Likelihood} & \textbf{Impact} & \textbf{Mitigation} \\
\midrule
$\psi$-correlation signal below detection threshold & Medium & High &
Increase $M$; improve single-shot sensitivity; extend to longer baselines \\
Environmental correlations mimic $\psi$ signal & Medium & High &
Extensive witness sensors; null tests; hardware exchange \\
Individual node failure & Low & Medium &
Modular design allows swap/repair; 4 nodes provides redundancy \\
Timing synchronization drift & Low & Low &
GPS common-view; periodic calibration against known tidal signals \\
DFD theory incorrect (no $\psi$ field) & Unknown & Existential &
Null result is scientifically valuable; apparatus can serve as
conventional gravity gradiometer network \\
\bottomrule
\end{tabular}
\caption{Risk assessment for the 4-node demonstrator.}
\label{tab:risk}
\end{table}


%=============================================================================
\section{Success Criteria}\label{sec:success}
%=============================================================================

\subsection{Minimum Success (Phase 1)}

\begin{enumerate}
\item Four nodes operational with single-shot phase noise $< 10^{-2}$ rad.
\item Environmental regression reduces classical correlations to
$< 10^{-6}$ level.
\item Tidal calibration agrees with theory to $< 10^{-7}$ fractional.
\item Null tests (time-reversal, hardware swap) pass at $5\sigma$.
\end{enumerate}

\subsection{Full Success (Phase 1)}

\begin{enumerate}
\item Excess inter-node correlation detected at $\geq 5\sigma$ significance.
\item Correlation follows DFD spatial dependence model
$C_{ij}(\mathbf{r}_{ij})$.
\item Correlation is independent of hardware identity (survives node swap).
\item Spectral shape consistent with $\psi$-field power spectrum prediction.
\end{enumerate}

\subsection{Extended Success (Phase 2/3)}

\begin{enumerate}
\item Sensitivity scaling verified: $N$-node sensitivity improves faster
than $N^{-1/2}$.
\item Gravitational wave sensitivity demonstrated in the 0.01--1 Hz band.
\item $\psi$-field tomography produces subsurface density maps consistent
with independent geological data.
\end{enumerate}


%=============================================================================
\section{Conclusion}\label{sec:conclusion}
%=============================================================================

This engineering design provides a complete, buildable specification for
a 4-node $\psi$-linked quantum sensor array demonstrator. The design uses
established cold-atom interferometry technology, requires no inter-node
coherent links, and is deployable in three phases of increasing baseline
and scientific ambition. The primary engineering challenge is not the
sensor technology itself---which is mature---but the systematic rejection
of classical correlation channels to the precision required to isolate
the $\psi$-field signal.

Whether the $\psi$-linked correlations predicted by Density Field Dynamics
are confirmed or ruled out, this apparatus represents a novel approach to
distributed quantum sensing that will produce valuable science at every
deployment phase: as a gravity gradiometer network for geophysics,
as a testbed for atom interferometer technology, and as a probe of
fundamental physics at the intersection of quantum mechanics and gravity.

\end{document}
